(1) 导体内部静电场必须为零。对金属内部 $a<r<b$ 的球面运用高斯定律，包围的总电量为 $q+Q_{\rm inner}=0$，故内表面有 $-q$，由壳净电荷 $Q$ 知外表面有 $Q+q$。球对称下腔内 $E_r=q/(4\pi\varepsilon_0r^2)$，金属中为零，外部 $E_r=(Q+q)/(4\pi\varepsilon_0r^2)$。
(2) 以无穷远电势为零，球壳电势 $V_{\rm shell}=(Q+q)/(4\pi\varepsilon_0b)$。因腔内由点电荷的电场积分得到 $V(r)-V(a)=q(1/r-1/a)/(4\pi\varepsilon_0)$，故 $V(r)=q(1/r-1/a)/(4\pi\varepsilon_0)+(Q+q)/(4\pi\varepsilon_0b)$。在 $r=a$ 与壳体常电势连接，没有人为引入电势跃变。
(3) 即使壳净电荷 $Q=0$，它的外表面也带 $+q$，外部仍有点电荷 $q$ 的库仑场，故中性不等于外部电场为零。球壳相对其净电荷的电势变化率 $dV_{\rm shell}/dQ=1/(4\pi\varepsilon_0b)$，电容 $C=4\pi\varepsilon_0b$。以 $Q=0$ 为参考，外部电源对送入电荷做功 $\Delta W=\int_0^QV_{\rm shell}(Q')\,dQ'=[Q^2/2+qQ]/(4\pi\varepsilon_0b)$。
(4) 理想点电荷自身场能 $\varepsilon_0\int E^2dV/2$ 在 $r\to0$ 发散，比较外壳位置或中心电荷位移时应先减去与位置无关的自能。中心电荷偏离后内壁表面电荷分布不再均匀，腔内场不再简单球对称；但高斯定律仍给内壁总量 $-q$、外壁总量 $Q+q$，球形等势外表面外场仍由总电量唯一决定。
